% page 141 of the facsimile (image p-140 of the scan)
% block 167.6 x 237.7 mm ; left margin 34.4 mm
\pgc{12.700mm}{0.000mm}{
\l{\cel{54}131}
\l{}
\l{}
\l{elongaco. Angulo-disto, mezurata de la centro di Tero, inter}
\l{\cel{3}Suno e ta o ca planeto.}
\l{}
\l{\sou{eloquenta}. I. Tre talentoza pri parolar. - II. (\sou{metaf.}) Qua}
\l{\cel{3}efikas quale parolo eloquenta. - DEFIS.}
\l{}
\l{\sou{eludar}. (\sou{trans.}) Evitar habile, rapide e ruze. - EFIS.}
\l{}
\l{\sou{\textquotedbl{}elzevir\textquotedbl{}}. (\sou{imprimarto}). Tipo gracioza, derivita de ta tipo}
\l{\cel{3}karakterizata da linei perpendikla (substitucita, en}
\l{\cel{3}1466, da du imprimisti Romana, a la tipi gotika). - DEFIRS.}
\l{}
\l{\sou{-em-}.\cel{2}Sufixo qua signifikas \textquotedbl{}qua tendencas a...\textquotedbl{}}
\l{}
\l{\sou{emalio}. I.Verniso qua rezultas de la vitrigo di substanci fuzebla}
\l{\cel{3}(silico, plomb-oxido, natroxo, potaso) quan onu aplikas, per}
\l{\cel{3}fuzo, sur la terakotajo, la fayencajo, la metalajo, ed a qua}
\l{\cel{3}onu adjuntas metal-oxidi pulverigita, por kolorizar lu diverse.}
\l{\cel{3}- II.\cel{2}Verniso formacita da la vitrigo ye la surfaco di}
\l{\cel{3}fayencajo. - DFIS.}
\l{}
\l{\sou{emanar}.(netrans., de). I. (aludante partikuli nepalpebla). Eska-}
\l{\cel{3}par de korpo sen ke la korpo percepteble deskreskas substan-}
\l{\cel{3}cale. - II. Devenar de ulo, de ulu. - III. (\sou{filoz}.)En la dok-}
\l{\cel{3}trino panteista pri emano, aludante la modi di substanco :}
\l{\cel{3}ekirar, per libereski intersucedanta, de la substanco una ed}
\l{\cel{3}universala qua esas Deo.-DEFIS.}
\l{}
\l{emancipar. (trans.) I. (legaro di la Roma antiqua).Liberigar}
\l{\cel{3}de la autoritato patrulala. - II. (\sou{legaro moderna}). Liberigar}
\l{\cel{3}(minoro) de la tutelo. - III. (en\cel{1}\sou{la Mezepoko}). Liberigar}
\l{\cel{3}(serfo, vasalo) de lua devi a lua sinioro. - IV. (\sou{metafore}).}
\l{\cel{3}Liberigar de la tutelo di autoritato supera. - DEFIRS.}
\l{}
\l{\sou{embalmar}. (\sou{trans.}) Prezervar (kadavro) de putro, per enpozar}
\l{\cel{3}substanci aromata, sikigiva ed antisepsiiva. - EF}
\l{}
\l{\sou{embarasar}. (\sou{trans., per, pro}).\cel{2}I. Jenar (pri agar, movar). -}
\l{\cel{3}II. Duktar a perplexeso, en qua lu ne savas quon decidor}
\l{\cel{3}o rezolvor. - EFIS.}
\l{}
\l{\sou{embargar}. (\sou{trans}.) I. (\sou{yuro marala}). Sorto di sequestro da}
\l{\cel{3}guvernantaro sur la navi di naciono stranjera kun qua lu mi-}
\l{\cel{3}litas. - II. (\sou{yuro civila}). Posedeskar, ye la nomo di la tri-}
\l{\cel{3}bunalo, mobli, imobli (di debanto por garantiar a la debario}
\l{\cel{3}to quo debesas a lu), objekti interdiktita, kontrabando-vari,}
\l{\cel{3}e c. - DEFIRS.}
\l{}
\l{\sou{embarkar}. I. (\sou{trans.}) Pozar en barko, en navo ; enirar navo por}
\l{\cel{3}voyajo sur aquo. (\sou{anke}\cel{1}: enirar treno, veturo). II. (\sou{netrans}.)}
\l{\cel{3}(\sou{aludante maraquo}.) Penetrar violentoze super la bordo di la}
\l{\cel{3}navo. - EFIS.}
}
